\uuid{DdCQ}
\exo7id{3838}
\titre{exo7 3838}
\auteur{quercia}
\organisation{exo7}
\datecreate{2010-03-11}
\isIndication{false}
\isCorrection{true}
\chapitre{Espace euclidien, espace normé}
\sousChapitre{Espaces vectoriels hermitiens}
\cpge{mpsi-mp2i.espaces-prehilbertiens.a, pcsi.espaces-prehilbertiens.a, ptsi.geometrie-plan.b, mpsi-mp2i.espaces-prehilbertiens.d, pcsi.espaces-prehilbertiens.d, mpsi-mp2i.changements-bases.a, pcsi.matrices-applications-lineaires.c, ptsi.matrices-applications-lineaires.c, mp-mpi.endomorphismes-euclidien.b, pc.endomorphismes-euclidiens.b, psi.endomorphismes-euclidien.b, pt.isometries-euclidien.b}
\module{Algèbre}
\niveau{L2}
\difficulte{}

\contenu{
\question{
Soit $E$ un espace euclidien, $(\vec e_1,\dots,\vec e_n)$ une base de $E$,
$G$ sa matrice de Gram et $G^{-1} = (a_{ij})$.

Montrer que :
$\forall\ \vec x\in E,\ \sum_{i,j} a_{ij}
  (\vec e_i\mid \vec x) (\vec e_j\mid \vec x) = \|\vec x\,\|^2$.
}
\reponse{
Soit $\cal B$ une base orthonormée de $E$ et $P$ la matrice de
passage de $\cal B$ à $(\vec e_i)$.

Le premier membre vaut ${}^tXPG^{-1}{}^tPX = {}^t XX$.
}
}
